% ============================================================
% boma.tex
% A Branch-Oriented Architecture for Mathematical Foundations:
% A Methodology for Constructing, Documenting, and Comparing
% Formal Theories
% ============================================================

\documentclass[12pt, a4paper]{article}

% ==============================
% PACKAGES
% ==============================
\usepackage{xcolor}
\usepackage[utf8]{inputenc}
\usepackage[T1]{fontenc}
\usepackage{lmodern}
\usepackage{microtype}

% Mathematics
\usepackage{amsmath}
\usepackage{amssymb}
\usepackage{amsthm}
\usepackage{mathtools}

% References and links
\usepackage[
  colorlinks=true,
  linkcolor=blue!60!black,
  citecolor=blue!60!black,
  urlcolor=blue!60!black
]{hyperref}
\usepackage{cleveref}

% Layout
\usepackage[
  top=2.5cm,
  bottom=2.5cm,
  left=3cm,
  right=3cm
]{geometry}
\usepackage{setspace}
\onehalfspacing

% Tables
\usepackage{booktabs}
\usepackage{array}

% Bibliography
\usepackage[
  backend=biber,
  style=numeric,
  sorting=nyt,
  maxbibnames=99
]{biblatex}
\addbibresource{boma.bib}

% ==============================
% THEOREM ENVIRONMENTS
% ==============================
\theoremstyle{plain}
\newtheorem{theorem}{Theorem}[section]
\newtheorem{proposition}[theorem]{Proposition}
\newtheorem{lemma}[theorem]{Lemma}
\newtheorem{corollary}[theorem]{Corollary}

\theoremstyle{definition}
\newtheorem{definition}[theorem]{Definition}
\newtheorem{axiom}[theorem]{Axiom}
\newtheorem{principle}[theorem]{Principle}

\theoremstyle{remark}
\newtheorem{remark}[theorem]{Remark}
\newtheorem{observation}[theorem]{Observation}
\newtheorem{application}[theorem]{Application}
\newtheorem{connection}[theorem]{Connection}

% ==============================
% CUSTOM COMMANDS
% ==============================
\newcommand{\FS}[1]{\mathfrak{F}(#1)}
\newcommand{\Br}[1]{\mathcal{B}_{#1}}
\newcommand{\Pa}[1]{\mathcal{P}_{#1}}
\newcommand{\DP}[1]{\mathcal{D}_{#1}}
\newcommand{\Th}[1]{\mathrm{Th}(#1)}
\newcommand{\Inj}[1]{\mathrm{Inj}(#1)}
\newcommand{\Der}[1]{\mathrm{Der}(#1)}
\newcommand{\IW}[1]{\omega(#1)}
\newcommand{\CIW}[1]{\Omega(#1)}
\newcommand{\DC}[1]{\overline{\mathrm{Dep}}(#1)}
\newcommand{\CC}[1]{\widehat{\mathrm{Content}}(#1)}
\newcommand{\CA}[2]{#1 \wedge #2}
\newcommand{\Sep}[1]{\mathrm{Sep}(#1)}
\newcommand{\BOMA}{\textsf{BOMA}}

% ==============================
% TITLE
% ==============================
\title{%
  \textbf{A Branch-Oriented Architecture for Mathematical Foundations}\\[0.5em]
  \large A Methodology for Constructing, Documenting,
  and Comparing Formal Theories%
}

\author{%
  [Author Name]\\
  \small [Institution]\\
  \small \texttt{[email]}%
}

\date{\today}

% ============================================================
\begin{document}
% ============================================================

\maketitle

% ==============================
% ABSTRACT
% ==============================
\begin{abstract}
Existing foundational systems for mathematics ---
axiomatic set theory, formal arithmetic, dependent type theories,
and categorical foundations --- are logically precise,
expressively powerful, and internally consistent.
What they do not provide, by design, is a structural record
of their own construction: the genuine choices made at critical
junctures of their development, the alternatives that were available
but not taken, and the specific assumptions that represent new
foundational commitments rather than derived consequences of prior
content. This absence is not a logical deficiency but an
architectural one.

We propose the \emph{Branch-Oriented Mathematical Architecture}
(\BOMA{}), a framework for the organization, documentation, and
comparison of mathematical foundations. \BOMA{} introduces seven
structural components: Blocks as the atomic units of foundational
content; a dependency relation partially ordering Blocks within a
development; Decision Points as explicit records of genuine
foundational choices; Branches as independent developments
originating from unchosen alternatives at Decision Points; Local
Consistency as the sole criterion of branch validity, assessed
independently of all other branches; Minimal Injection as the
discipline governing the introduction of new assumptions; and
Architectural Transparency as the conjunctive structural property
produced by the full application of these requirements.

A central consequence of the framework is the Foundational Space:
the collection of all locally consistent architectural paths
sharing a common origin, equipped with a natural
ultra-pseudometric structure. Within this space, contradictions
between branches are not architectural failures but structural
information --- precise measurements of the non-triviality of the
Decision Points at which diverging branches separate.
Comparison between branches is established as an epistemological
instrument, not a validation procedure.

Every existing foundational system is shown to correspond to a
locally consistent, architecturally opaque path within an
appropriate Foundational Space. \BOMA{} is a meta-framework: it
makes no specific mathematical commitments and competes with no
existing foundation. It adds a structural dimension that the
existing paradigm does not supply, and constitutes a possibility
space within which unknown mathematics occupies a determinate,
reachable, and structurally characterized position.
\end{abstract}

\medskip
\noindent\textbf{MSC 2020:}
03A05, 03B22, 03B70, 00A30, 03F65.

\medskip
\noindent\textbf{Keywords:}
mathematical foundations,
branch-oriented architecture,
decision points,
local consistency,
foundational space,
architectural transparency,
minimal injection,
\BOMA{}.

\tableofcontents
\newpage

% ============================================================
\section{Introduction}
\label{sec:introduction}
% ============================================================

\subsection{The Success and the Silence}
\label{subsec:success}

The foundational systems of modern mathematics are among the
most carefully constructed intellectual artifacts in the history
of human thought. Zermelo-Fraenkel set theory with the Axiom of
Choice~\cite{Jech2003,Zermelo1908}, the Peano axioms in their
first-order formulation~\cite{Peano1889}, Martin-L\"{o}f Type
Theory~\cite{MartinLof1975,MartinLof1984}, Homotopy Type
Theory~\cite{HoTTBook2013}, and the categorical foundations
developed from Lawvere's work~\cite{Lawvere1964,MacLane1971} ---
each of these systems is internally precise to a degree that
admits mechanical verification, expressively powerful enough to
formalize the vast majority of contemporary mathematical practice,
and consistent by all available evidence.

This success deserves acknowledgment before anything else is said.
The architecture proposed in this paper is addressed to
researchers who work within, or in close proximity to, this
foundational tradition. Its starting point is not a failure of
existing foundations but an observation about what those
foundations, by the nature of their design, do not provide.

The observation is this: \emph{existing foundational systems are
silent about their own construction.}

A mathematician who opens a standard presentation of ZFC
encounters nine axioms, stated in the language of first-order
logic with a single primitive relation. The axioms are clearly
labeled. The inference rules are declared. The derivations are
checkable. Nothing in this presentation indicates that the
adoption of the Axiom of Choice was contested, that the Axiom of
Foundation was introduced to exclude pathological sets, that the
first-order formulation was selected over a second-order
alternative, or that the underlying logic could have been chosen
to be intuitionistic rather than classical. These choices were
made. Their consequences permeate every theorem derivable within
ZFC. But the presentation of ZFC does not record them.

\subsection{What Is Not Visible}
\label{subsec:invisible}

The silence identified in \S\ref{subsec:success} has three
specific manifestations.

\textbf{The invisibility of Decision Points.}
Every foundational system passed through moments of genuine
choice. These moments --- which we term \emph{Decision Points}
--- are real architectural events. They determined the character
of the mathematics that followed. But in the formal presentations
of existing systems, they leave no trace.

\textbf{The suppression of alternatives.}
A consequence of invisible Decision Points is that the space of
foundational alternatives surrounding any existing system is not
visible from within that system's formal presentation.

\textbf{The conflation of injections and derivations.}
Existing foundational presentations do not systematically
distinguish between statements that are genuine new assumptions
and statements that are derived consequences of earlier
assumptions.

\begin{remark}[Prior Work on Theory Structuring]
\label{rem:prior-work}
Several existing research programs address aspects of the
problem identified above.
Theory structuring languages --- most prominently
\textsc{casl}~\cite{SannellaTarlecki2012} and its predecessors
--- provide formal mechanisms for building theories from smaller
components. The \textsc{omdoc} format and the
\textsc{mmt} framework~\cite{Kohlhase2006,RabeKohlhase2013}
provide comprehensive approaches to the representation of
mathematical knowledge including theory graphs.
Institution theory~\cite{GoguenBurstall1992} provides a
categorical framework for logical systems achieving a form of
indifference to the specific logic employed.

The present paper does not claim that these frameworks are
superseded by \BOMA{}. Each addresses a genuine problem in the
organization of formal mathematical knowledge. \BOMA{} addresses
a problem that none of them addresses: the explicit, structurally
mandatory documentation of foundational decisions within the
formal record of a mathematical development, together with the
preservation of unchosen alternatives as first-class
architectural objects.
\end{remark}

\subsection{The Architectural Question}
\label{subsec:architecture}

The three absences described above suggest a question that is
not, in the first instance, a mathematical question:
\begin{quote}
\emph{Is the linear presentation of foundational mathematics the
only organizational paradigm consistent with mathematical rigor?
Or is there an alternative paradigm that preserves all the
logical properties of the existing approach while adding a
structural dimension that the existing approach does not supply?}
\end{quote}
This is an architectural question. The present paper proposes
such a framework: the Branch-Oriented Mathematical Architecture,
hereafter \BOMA{}.

The answer to this question has consequences beyond the
organization of existing mathematics. A framework that makes
foundational construction transparent and branchable is
simultaneously a framework for mathematical exploration: every
documented Decision Point is an entrance to mathematical
territory that may not yet have been traversed.

\subsection{Our Contribution}
\label{subsec:contribution}

This paper makes the following contributions.

\textbf{A formal definition of Architectural Transparency.}
We introduce a precise characterization of what it means for a
foundational development to be fully legible: to expose its
dependency structure, its Decision Points, its injection
structure, and its Branch Tree in a form computable from the
structural record without reconstruction.

\textbf{A structural framework for foundational development.}
We introduce seven structural components --- Blocks, Dependencies,
Decision Points, Branches, Local Consistency, Minimal Injection,
and Architectural Transparency --- and develop their formal
properties in \S\ref{sec:architecture}.

\textbf{The principle of Local Consistency as the sole criterion
of branch validity.} We establish in \S\ref{subsec:local} that a
branch is valid if and only if it is internally consistent,
independently of all other branches.

\textbf{The Foundational Space as a mathematical object.}
We introduce in \S\ref{sec:branching} the Foundational Space
$\FS{\beta_0}$: the collection of all locally consistent
architectural paths sharing a common Atomic Block, equipped with
a natural ultra-pseudometric structure.

\textbf{A precise account of comparison as an epistemological
tool.} We establish in \S\ref{sec:comparison} that structural
comparison between branches is not a validation procedure but an
epistemological instrument.

\textbf{The situation of existing foundations within \BOMA{}.}
We show in \S\ref{sec:existing} that every existing foundational
system corresponds to a locally consistent, architecturally opaque
path in an appropriate Foundational Space.

A three-level logical architecture --- object logic, framework
logic, and meta-logic --- is introduced in
Remark~\ref{rem:three-levels}, establishing the precise scope of
the framework's logical neutrality claims.

\subsection{What This Paper Does Not Claim}
\label{subsec:notclaim}

This paper does not claim that existing foundations are
incorrect, incomplete, or in need of correction. It does not
claim that \BOMA{} is superior to any existing foundational
system. It does not claim to resolve the philosophical questions
surrounding foundational pluralism. It does not claim that the
applications surveyed in \S\ref{sec:applications} have been
realized. The claim is more modest: that there is an
architectural dimension of foundational mathematics that has been
systematically invisible within the existing paradigm, that this
dimension can be made precise, and that a framework for doing so
can be developed with sufficient rigor.

\subsection{Paper Organization}
\label{subsec:organization}

Section~\ref{sec:paradigm} characterizes the existing
foundational paradigm. Section~\ref{sec:architecture} develops
\BOMA{} in full. Section~\ref{sec:branching} introduces the
Foundational Space. Section~\ref{sec:comparison} develops the
theory of structural comparison. Section~\ref{sec:existing}
situates existing foundations within \BOMA{}.
Section~\ref{sec:applications} surveys applications and research
programs. Section~\ref{sec:conclusion} concludes.

Several cross-cutting remarks developed in the course of the
paper address issues that span multiple sections.
Remark~\ref{rem:three-levels} introduces a three-level logical
architecture. Remark~\ref{rem:reconstruction} introduces
architectural reconstruction as the interpretive methodology
governing the analyses of \S\ref{sec:existing}.
Remark~\ref{rem:possibility} establishes the Foundational Space
as a possibility space. Remark~\ref{rem:scope} articulates the
full scope of the framework's vision. These remarks are integral
to the paper's argument.

% ============================================================
\section{The Current Paradigm}
\label{sec:paradigm}
% ============================================================

\subsection{The Structure of Existing Foundations}
\label{subsec:structure}

\begin{observation}[The Linear Presentation]
\label{obs:linear}
Every major foundational system is presented as a
\emph{linear development}~\cite{Enderton2001,Shoenfield1967}:
a sequence of declarations ordered so that each element
presupposes only what has been declared before it. The sequence
begins at a fixed starting point and extends forward without
branching.

A compatible linearization of the Block dependency structure ---
as defined in Definition~\ref{def:linearization} --- is precisely
what a traditional textbook presentation provides: one admissible
reading order among many. Different textbooks on the same
foundational system may select different linearizations, but all
are compatible with the same underlying dependency structure.
\end{observation}

\begin{observation}[Primitive Notions and Axioms as Fixed Points]
\label{obs:fixed}
In every existing foundational system, the primitive notions and
initial axioms are presented as fixed points: given without
derivation and treated as the terminal points of any regress of
justification.
\end{observation}

\subsection{What Existing Foundations Achieve}
\label{subsec:achieve}

\begin{observation}[Consistency and Coherence]
\label{obs:consistency}
The primary achievement of existing foundational systems is the
provision of a framework within which mathematical reasoning can
be conducted without risk of contradiction~\cite{Goedel1931,Goedel1938}.
\end{observation}

\begin{observation}[Expressive Power]
\label{obs:expressive}
Existing foundational systems are expressively powerful to an
extent that continues to exceed the demands placed on them by
mathematical practice~\cite{Jech2003,MartinLof1984}.
\end{observation}

\begin{observation}[Computability and Mechanization]
\label{obs:mechanization}
The formal specification of existing foundations has enabled the
development of proof assistants.
Coq~\cite{Coq2023}, Lean~\cite{deMouraEtAl2015},
Agda~\cite{Norell2007}, Isabelle~\cite{NipkowEtAl2002},
and Metamath~\cite{Megill1997} are among the systems that have
made this possible.
\end{observation}

\begin{observation}[Theoretical Comparability via Interpretability]
\label{obs:interpretability}
Existing foundational systems have been compared through the
notion of relative interpretability~\cite{Shoenfield1967,Kunen1980}.
\end{observation}

\subsection{The Architectural Gap}
\label{subsec:gap}

\begin{observation}[The Invisibility of Decision Points]
\label{obs:invisibility}
Every existing foundational system passed through moments of
genuine choice. These moments are not invented by the present
framework. They are episodes in the history of mathematics,
documented in primary sources~\cite{Moore1982,vanHeijenoort1967,Zermelo1904}.
What they share is a common structure: a moment at which more
than one mathematically coherent path was available, a process
of deliberation, and a selection. We term this structure a
\emph{Decision Point}.

The observation that motivates the present paper is not that
existing formal systems fail to contain a concept we have
introduced. It is that the formal presentations of existing
systems provide no structural mechanism for recording, within
the formalism itself, the Decision Points that their historical
development passed through.
\end{observation}

\begin{observation}[The Suppression of Alternatives]
\label{obs:suppression}
A consequence of Observation~\ref{obs:invisibility} is that the
alternatives to the choices actually made are also invisible.
\end{observation}

\begin{observation}[The Absence of Architectural History]
\label{obs:history}
The combination of invisible Decision Points and suppressed
alternatives produces the \emph{absence of architectural history}:
the phenomenon in which a foundational system presents its
current state without any record of how it came to be in that
state.
\end{observation}

\begin{observation}[The Conflation of Injections and Derivations]
\label{obs:conflation}
Existing foundational presentations do not distinguish between
the assumptions they inject and the theorems they derive in terms
of their architectural status.
\end{observation}

\subsection{The Nature of the Gap: Architectural, Not Logical}
\label{subsec:nature}

\begin{remark}
The gap identified in \S\ref{subsec:gap} is not a logical gap.
Existing foundational systems are logically complete in the sense
relevant to their purposes. The gap is \emph{architectural}:
existing foundations are organized according to a paradigm that
optimizes for logical consistency, expressive power, and
mechanizability --- but not for transparency of Decision Points,
preservation of alternatives, or explicit injection structure.
\end{remark}

\begin{remark}
The architectural gap is not a defect in existing foundations;
it is a dimension in which they are silent. The architecture
proposed in this paper does not correct existing foundations;
it extends the paradigm within which foundations are organized.
\end{remark}

\subsection{Existing Foundations Within the Architectural Paradigm}
\label{subsec:existing-in-boma}

\begin{observation}
\label{obs:existing-paths}
Every existing foundational system corresponds, within the
architectural framework, to a locally consistent path in some
Foundational Space $\FS{\beta_0}$.
\end{observation}

\begin{observation}
\label{obs:existing-opaque}
Existing foundational systems correspond to
\emph{architecturally opaque} paths: paths that satisfy local
consistency but do not satisfy Architectural Transparency,
because their Decision Points are not documented and their Branch
Trees are not recorded.

The application of the architectural framework to existing
foundations is an instance of what \S\ref{sec:existing} terms
\emph{architectural reconstruction}: a retrospective
interpretive reading, not a mechanical extraction.
\end{observation}

% ============================================================
\section{The Proposed Architecture (\BOMA{})}
\label{sec:architecture}
% ============================================================

\subsection{Blocks}
\label{subsec:blocks}

\subsubsection*{Motivation}

Every formal development of mathematics organizes its content
into primitive notions, axioms, definitions, theorems, and
inference rules. What is systematically absent is an explicit
account of the \emph{unit of packaging} --- the structure that
groups related content together, records what that content
depends upon, and tracks whether it is stable or subject to
revision. That structure is the \emph{Block}.

\begin{definition}[Block]
\label{def:block}
A \emph{Block} is a four-tuple:
\[
  \beta
  = \bigl(\,\mathrm{Id}(\beta),\;
           \mathrm{Content}(\beta),\;
           \mathrm{Dep}(\beta),\;
           \mathrm{Status}(\beta)\,\bigr)
\]
where:
\begin{itemize}
  \item $\mathrm{Id}(\beta)$ is a unique identifier from a
        totally ordered index set $\mathcal{I}$;
  \item $\mathrm{Content}(\beta)$ is a finite, non-empty
        collection of formal statements in a declared language
        $\mathcal{L}$;
  \item $\mathrm{Dep}(\beta) \subseteq
        \{\beta' \mid \mathrm{Id}(\beta') < \mathrm{Id}(\beta)\}$
        is the set of Blocks directly presupposed by
        $\mathrm{Content}(\beta)$;
  \item $\mathrm{Status}(\beta) \in \{\mathrm{Fixed},\,
        \mathrm{Open}\}$.
\end{itemize}
\end{definition}

\begin{definition}[Fixed Block]
\label{def:fixed}
A Block $\beta$ is \emph{Fixed} if its content is stable within
the current architectural path. A Fixed Block is a reliable
node; any Block depending on it may do so without anticipating
changes.
\end{definition}

\begin{definition}[Open Block]
\label{def:open}
A Block $\beta$ is \emph{Open} if its content is subject to
revision. An Open Block may transition to Fixed once its content
is stabilized. The reverse transition is not permitted within a
given path; revision is effected by branching, not by modifying
in place.
\end{definition}

\begin{definition}[Atomic Block]
\label{def:atomic}
A Block $\beta$ is \emph{Atomic} if $\mathrm{Dep}(\beta) =
\emptyset$. Atomic Blocks are the roots of the architectural
development, carrying primitive notions, initial axioms, and the
declaration of the underlying logic.
\end{definition}

\begin{remark}
\label{rem:atomic-significance}
Every architectural development must contain at least one
Atomic Block. The Atomic Block is the location at which the most
fundamental choices are made explicit: what counts as a primitive
notion, what logic governs inference, what is accepted without
derivation. In traditional foundational presentations, these
choices are made implicitly. In \BOMA{}, they are the declared
content of the first Block.
\end{remark}

\begin{definition}[Block Identity and Content-Equivalence]
\label{def:identity}
Two Blocks $\beta$ and $\beta'$ are \emph{identical} if and only
if $\mathrm{Id}(\beta) = \mathrm{Id}(\beta')$. They are
\emph{content-equivalent} if $\mathrm{Content}(\beta) =
\mathrm{Content}(\beta')$ up to notational equivalence.
\end{definition}

\begin{remark}[Blocks and Theory Structuring]
\label{rem:theory-structuring}
The notion of a Block as defined in
Definition~\ref{def:block} is related to, but distinct from,
the notion of a structured theory presentation in algebraic
specification languages such as
\textsc{casl}~\cite{SannellaTarlecki2012} and in mathematical
knowledge management systems such as
\textsc{omdoc}~\cite{Kohlhase2006} and
\textsc{mmt}~\cite{RabeKohlhase2013}.
In these systems, a theory extension records what was added;
a Block in \BOMA{} records, additionally, the injection structure
of the added content (\S\ref{subsec:injection}) and --- when the
Block is a Decision Block --- the alternatives that were not
chosen and the rationale for the selection
(\S\ref{subsec:decision}).
\end{remark}

\subsection{Dependencies and Ordering}
\label{subsec:dependencies}

\begin{definition}[Direct Dependency]
\label{def:direct-dep}
$\beta'$ \emph{directly depends on} $\beta$, written
$\beta \prec_1 \beta'$, if and only if
$\beta \in \mathrm{Dep}(\beta')$.
\end{definition}

\begin{definition}[Dependency Relation]
\label{def:dep-relation}
The \emph{dependency relation} $\prec$ is the transitive closure
of $\prec_1$. We say $\beta$ is an \emph{ancestor} of $\beta'$
and $\beta'$ is a \emph{descendant} of $\beta$.
\end{definition}

\begin{definition}[Dependency Closure]
\label{def:dep-closure}
The \emph{dependency closure} of $\beta$ is:
\[
  \DC{\beta}
  = \{\,\beta' \mid \beta' \prec \beta\,\}.
\]
\end{definition}

\begin{axiom}[Acyclicity]
\label{ax:acyclic}
The dependency relation $\prec$ is irreflexive and acyclic.
\end{axiom}

\begin{proposition}
\label{prop:wellfounded}
If $\prec$ is acyclic on a finite set of Blocks, then it is
well-founded: every non-empty subset contains at least one Block
with no ancestors within that subset. The minimal elements are
precisely the Atomic Blocks of
Definition~\ref{def:atomic}.
\end{proposition}

\begin{proof}
Let $\{\beta_0,\ldots,\beta_n\}$ be a finite set of Blocks with
an acyclic dependency relation $\prec$. Consider any non-empty
subset $S$. Since $S$ is finite and $\prec$ is acyclic on $S$,
there are no infinite descending chains. A finite acyclic
relation is well-founded, so every non-empty finite set with an
irreflexive transitive relation contains at least one element
with no predecessors within that set. Such $\beta^* \in S$
satisfies: there is no $\beta' \in S$ with $\beta' \prec
\beta^*$. Taking $S$ to be the full set of Blocks, the minimal
elements are those with no predecessors in the entire path,
i.e., Blocks $\beta$ for which $\mathrm{Dep}(\beta) = \emptyset$.
These are precisely the Atomic Blocks.
\end{proof}

\begin{proposition}
\label{prop:partial-order}
The dependency relation $\prec$ together with the identity
relation defines a strict partial order on the Blocks of any
architectural path.
\end{proposition}

\begin{proof}
Irreflexivity follows from Axiom~\ref{ax:acyclic}. Transitivity
holds by construction. For asymmetry: suppose $\beta \prec \beta'$
and $\beta' \prec \beta$. By transitivity of $\prec$, we obtain
$\beta \prec \beta$, contradicting irreflexivity. Hence $\prec$
is asymmetric, and together with irreflexivity and transitivity,
constitutes a strict partial order.
\end{proof}

\begin{remark}
The structure is a \emph{partial} order, not a total one. In a
partial order, several Blocks may be mutually independent and
developed in parallel. This reflects the actual structure of
mathematical knowledge more faithfully than a linear presentation.
\end{remark}

\begin{definition}[Compatible Linearization]
\label{def:linearization}
A \emph{compatible linearization} of the Blocks of a path is a
total order $<_\ell$ such that
$\beta \prec \beta' \implies \beta <_\ell \beta'$.
Compatible linearizations exist for every finite acyclic
dependency structure (standard topological sorting
theorem~\cite{Enderton2001}).
\end{definition}

\begin{definition}[Cumulative Content]
\label{def:cumulative}
The \emph{cumulative content} of $\beta$ is:
\[
  \CC{\beta}
  = \mathrm{Content}(\beta)
  \;\cup
  \bigcup_{\beta' \in \DC{\beta}} \mathrm{Content}(\beta').
\]
\end{definition}

\subsection{Decision Points}
\label{subsec:decision}

\subsubsection*{Motivation}

Every foundational development contains moments of genuine
choice. The architecture requires that every genuine choice be
made explicit, recorded at the moment it is made, and preserved
as a permanent feature of the architectural path.

\begin{definition}[Decision Point]
\label{def:decision-point}
A \emph{Decision Point} is a five-tuple:
\[
  \mathcal{D}
  = \bigl(\,\mathrm{Id}(\mathcal{D}),\;
           \mathrm{Anchor}(\mathcal{D}),\;
           \mathrm{Options}(\mathcal{D}),\;
           \mathrm{Selection}(\mathcal{D}),\;
           \mathrm{Rationale}(\mathcal{D})\,\bigr)
\]
where $|\mathrm{Options}(\mathcal{D})| \geq 2$.
\end{definition}

\begin{remark}
The $\mathrm{Rationale}$ component is not required to be a formal
proof that the selection is superior to the alternatives. It is a
transparent account of the considerations that guided the
selection --- mathematical, pragmatic, historical, or a
combination.
\end{remark}

\begin{axiom}[Mandatory Documentation]
\label{ax:documentation}
Every Decision Point must be fully documented at the moment it
is introduced:
\begin{enumerate}
  \item $\mathrm{Options}(\mathcal{D})$ enumerates all
        genuinely considered alternatives;
  \item $\mathrm{Selection}(\mathcal{D})$ is unambiguous;
  \item $\mathrm{Rationale}(\mathcal{D})$ is explicit;
  \item Unchosen options are permanently preserved.
\end{enumerate}
\end{axiom}

\begin{definition}[Forced Step and Genuine Choice]
\label{def:genuine-choice}
A step from Block $\beta$ is \emph{forced} if the content of
the next Block is uniquely determined by prior cumulative content.
A step is a \emph{genuine choice} --- objectively, if the
alternatives' logical admissibility and non-equivalence is
established by a formal independence result; \emph{interpretively},
if the developer judges it so without a formal independence proof.
\end{definition}

\begin{remark}[The Interpretive Character of Decision Point Identification]
\label{rem:interpretive}
The distinction between forced steps and genuine choices is not
always decidable from within the formal record of the path. Both
objectively and interpretively identified Decision Points are
legitimate, but they carry different epistemic weights. The basis
of identification must be declared in the
$\mathrm{Rationale}$ component of the Decision Block.

This interpretive dimension is not a deficiency of the framework.
It reflects an irreducible feature of foundational construction:
the assessment of what is genuinely optional is itself a
foundational act. The framework does not eliminate interpretive
judgment; it makes it visible and auditable.
\end{remark}

\begin{definition}[Decision Block]
\label{def:decision-block}
A \emph{Decision Block} is a Block whose content is precisely the
documentation of a Decision Point satisfying
Axiom~\ref{ax:documentation}.
\end{definition}

\begin{definition}[Decision Point Catalogue]
\label{def:catalogue}
The \emph{Decision Point Catalogue} of an architectural path is
the ordered collection of all Decision Blocks, ordered by
position in the dependency structure. It constitutes the
complete record of every genuine choice made in the development.
\end{definition}

\begin{remark}
In existing foundational systems, the analogue of a Decision
Point Catalogue does not exist as an explicit artifact. The
choices that shaped ZFC, Peano Arithmetic, or the Calculus of
Constructions are recoverable only through historical scholarship.
\BOMA{} makes the catalogue a first-class structural object.
\end{remark}

\subsection{Branches}
\label{subsec:branches}

\begin{definition}[Path]
\label{def:path}
An \emph{architectural path} $\mathcal{P}$ is a finite sequence
of Blocks $(\beta_0, \beta_1, \ldots, \beta_n)$ satisfying:
\begin{enumerate}
  \item $\beta_0$ is Atomic;
  \item For each $i > 0$, $\mathrm{Dep}(\beta_i) \subseteq
        \{\beta_0, \ldots, \beta_{i-1}\}$;
  \item The dependency relation is acyclic
        (Axiom~\ref{ax:acyclic});
  \item Every Decision Block satisfies
        Axiom~\ref{ax:documentation}.
\end{enumerate}
The local consistency requirement referenced here is defined
formally in \S\ref{subsec:local}. This forward dependency is
unavoidable given that the validity criterion of
\S\ref{subsec:local} presupposes the path structure of this
section; the two sections are mutually dependent in their
motivation and should be read as a unified specification.
\end{definition}

\begin{definition}[Trunk Path]
\label{def:trunk}
The \emph{Trunk Path} $\mathcal{P}_0$ is the initial
architectural path --- the path produced by following, at every
Decision Point, the option recorded as
$\mathrm{Selection}(\mathcal{D})$.
\end{definition}

\begin{definition}[Branching Prefix]
\label{def:prefix}
Let $\mathcal{P}$ contain a Decision Block $\beta_j$ for
$\mathcal{D}$. The \emph{branching prefix} of $\mathcal{P}$ at
$\mathcal{D}$ is the initial segment
$\mathcal{P}_{|\mathcal{D}} = (\beta_0, \ldots, \beta_j)$.
\end{definition}

\begin{definition}[Branch]
\label{def:branch}
Let $\mathcal{P}$ contain Decision Point $\mathcal{D}$ with
options $\{O_1,\ldots,O_k\}$ and selection $O_s$. A
\emph{Branch} at $\mathcal{D}$ selecting option $O_i \neq O_s$
is an architectural path:
\[
  \mathcal{B}_i
  = (\beta_0, \ldots, \beta_j,\,
     \beta^i_{j+1}, \beta^i_{j+2}, \ldots)
\]
where $(\beta_0,\ldots,\beta_j) = \mathcal{P}_{|\mathcal{D}}$
is the branching prefix inherited exactly, and subsequent Blocks
are developed under the commitment of $O_i$.
\end{definition}

\begin{definition}[Branch at Depth $d$]
\label{def:branch-depth}
A Branch at the $m$-th Decision Point (ordered by position in
the dependency structure) is said to be at \emph{depth} $m$.
Branches at greater depth share a longer common ancestry and
inherit a more developed foundational context before diverging.
\end{definition}

\begin{proposition}[Inheritance of Branching Prefix]
\label{prop:inheritance}
Let $\mathcal{B}_i$ be a Branch at $\mathcal{D}$ with branching
prefix $(\beta_0,\ldots,\beta_j)$. Then every Block $\beta_\ell$
for $\ell \leq j$ appears in $\mathcal{B}_i$ with the same
identifier, content, dependency set, and status as in
$\mathcal{P}$; and $\CC{\beta_j}$ is identical in both.
\end{proposition}

\begin{proof}
Immediate from Definition~\ref{def:branch}: the branching prefix
is inherited exactly, with no modification to any Block's
components.
\end{proof}

\begin{principle}[Branch Independence]
\label{prin:independence}
The development of $\mathcal{B}_i$ after $\mathcal{D}$ is fully
independent of $\mathcal{P}$ after $\mathcal{D}$: no post-Decision
Block of $\mathcal{P}$ is available to $\mathcal{B}_i$ unless
independently derived within $\mathcal{B}_i$; and no result of
$\mathcal{B}_i$ affects or constrains $\mathcal{P}$.
\end{principle}

\begin{definition}[Branch Tree]
\label{def:branch-tree}
The \emph{Branch Tree} $\mathcal{T}$ of a development is the
rooted tree whose root is the Atomic Block $\beta_0$, whose
internal nodes are Decision Blocks, whose edges are labeled by
the options $O_i \in \mathrm{Options}(\mathcal{D})$, and whose
paths from root to any node correspond to unique architectural
paths.
\end{definition}

\begin{remark}
The Foundational Space $\FS{\beta_0}$ defined formally in
\S\ref{sec:branching} is the collection of all locally consistent
paths sharing $\beta_0$. The Branch Tree is the architectural
record of one development within that space.
\end{remark}

\begin{definition}[Minimal Branch]
\label{def:minimal-branch}
A Branch $\mathcal{B}_i$ at $\mathcal{D}$ is \emph{minimal} if
the first Post-Decision Block differs from the corresponding
Block in the parent path only in the content directly determined
by the selection of $O_i$, and all subsequent Blocks are
developed under no additional assumptions beyond the branching
prefix and the commitment of $O_i$.
\end{definition}

\subsection{Local Consistency as the Sole Validity Criterion}
\label{subsec:local}

\begin{remark}[Three Logical Levels]
\label{rem:three-levels}
The architecture operates simultaneously at three logical levels.

The \emph{object logic} is declared within each individual path
at its Atomic Block. The architecture is genuinely indifferent
to the choice of object logic.

The \emph{framework logic} is the logic used to define the
architectural concepts themselves. The present paper employs
classical first-order logic~\cite{Enderton2001,Shoenfield1967}
as its framework logic. This is a choice, not a necessity: a
version of the framework formulated in intuitionistic logic would
be a coherent alternative, constituting a Branch from the present
formulation at the Decision Point of framework logic selection.

The \emph{meta-logic} is the logic in which the propositions and
proofs of this paper are expressed; it is likewise classical.

The claim of logical neutrality applies exclusively to the object
logic level.
\end{remark}

\begin{definition}[Deductive Closure of a Branch]
\label{def:closure}
The \emph{deductive closure} of $\mathcal{B}$ is:
\[
  \Th{\mathcal{B}}
  = \{\,\varphi \mid \mathcal{B} \vdash \varphi\,\}
\]
where $\vdash$ denotes derivability under the object logic
declared in the Atomic Block of $\mathcal{B}$.
\end{definition}

\begin{definition}[Local Consistency]
\label{def:local-consistency}
A branch $\mathcal{B}$ is \emph{locally consistent} if there is
no formula $\varphi$ in the language of the path's declared
object logic such that both $\mathcal{B} \vdash \varphi$ and
$\mathcal{B} \vdash \neg\varphi$, where $\neg$ denotes negation
as defined within that object logic.
\end{definition}

\begin{remark}
The use of $\neg\varphi$ in Definition~\ref{def:local-consistency}
is relative to the object logic of the path. For logics without a
standard negation, an appropriate analogue must be declared in
the Atomic Block.
\end{remark}

\begin{definition}[Branch Validity]
\label{def:validity}
A branch $\mathcal{B}$ is \emph{valid} if and only if it is
locally consistent. No further condition is imposed.
\end{definition}

\begin{remark}[On the Use of ``Principle'']
\label{rem:principle}
Throughout this paper, the term \emph{Principle} designates an
architectural commitment that is neither a primitive assumption
(Axiom) nor a derived result (Proposition), but a structural
norm: a requirement that follows from the design goals of the
framework and governs how its components are to be used.
\end{remark}

\begin{remark}[Minimality of the Validity Criterion]
Branch validity requires neither consistency with any external
system, nor agreement with any other branch, nor completeness.
Local consistency is both necessary and sufficient.
\end{remark}

\begin{principle}[Independence of Validity]
\label{prin:validity-independent}
The local consistency of $\mathcal{B}_1$ is logically independent
of the local consistency of $\mathcal{B}_2$. If
$\mathcal{B}_1 \vdash \varphi$ and
$\mathcal{B}_2 \vdash \neg\varphi$, both branches may nonetheless
be locally consistent, and therefore both valid.
\end{principle}

\begin{principle}[Contradiction as Architectural Information]
\label{prin:contradiction}
Let $\mathcal{B}_1$ and $\mathcal{B}_2$ be valid branches with:
\[
  \mathcal{B}_1 \vdash \varphi
  \qquad\text{and}\qquad
  \mathcal{B}_2 \vdash \neg\varphi.
\]
This situation is not a failure of the architecture. It is the
precise content the architecture is designed to surface: the
Decision Point $\mathcal{D}$ at which they diverged was
\emph{non-trivial}. Both mathematical worlds are coherent on
their own terms, and the architecture preserves both without
forcing a resolution.
\end{principle}

\begin{definition}[Non-trivial Decision Point]
\label{def:nontrivial}
A Decision Point $\mathcal{D}$ is \emph{non-trivial} if there
exist two valid branches $\mathcal{B}_1$ and $\mathcal{B}_2$
diverging at $\mathcal{D}$ such that for some formula $\varphi$:
$\mathcal{B}_1 \vdash \varphi$ and
$\mathcal{B}_2 \vdash \neg\varphi$.
\end{definition}

\begin{remark}
Comparison between branches is addressed to two or more branches:
what do their differences reveal? This is an epistemological
question, not a logical one. Comparison is a tool for learning,
not for adjudication. There is no mechanism by which one branch
can invalidate another.
\end{remark}

\subsection{Minimal Injection}
\label{subsec:injection}

\begin{definition}[Injection]
\label{def:injection}
A statement $\varphi \in \mathrm{Content}(\beta)$ is an
\emph{injection} of $\beta$ if:
\[
  \CC{\DC{\beta}} \nvdash \varphi.
\]
We write:
\[
  \Inj{\beta}
  = \bigl\{\,\varphi \in \mathrm{Content}(\beta)
    \mid \CC{\DC{\beta}} \nvdash \varphi\,\bigr\},
  \qquad
  \Der{\beta} = \mathrm{Content}(\beta) \setminus \Inj{\beta}.
\]
\end{definition}

\begin{definition}[Injection Weight]
\label{def:injection-weight}
The \emph{injection weight} of $\beta$ is
$\IW{\beta} = |\Inj{\beta}|$.
\end{definition}

\begin{axiom}[Minimal Injection Principle]
\label{ax:minimal}
Every Block $\beta$ must satisfy:
\begin{enumerate}
  \item \textbf{Necessity}: every injection is necessary for the
        stated purpose of $\beta$;
  \item \textbf{Individuation}: every injection is declared
        individually and explicitly;
  \item \textbf{Non-derivability Acknowledgment}: the fact that
        $\varphi$ is an injection is acknowledged;
  \item \textbf{Minimality}: if $\varphi$ can be replaced by a
        logically weaker $\psi$ sufficient for the purpose of
        $\beta$, then $\psi$ is injected instead.
\end{enumerate}
\end{axiom}

\begin{remark}[Injection Weight and Computability]
\label{rem:computability}
The injection set $\Inj{\beta}$ is well-defined but subject to
an important qualification. In the general case of first-order
logic, determining whether $\varphi$ is derivable from a set of
prior statements is undecidable, as an immediate consequence of
Church's theorem~\cite{Church1936}. The injection set is
therefore well-defined but not mechanically computable in full
generality.

This qualification does not undermine the architectural
significance of injection weights. In many practically relevant
cases --- finitely axiomatized systems with decidable consequence
relations --- injection weights are effectively computable. The
value of the Minimal Injection Principle lies in the
architectural discipline it enforces: the developer is required
to identify and declare each injection individually. This
discipline is independent of whether a mechanical verifier could
confirm the non-derivability claim.
\end{remark}

\begin{definition}[Cumulative Injection Weight]
\label{def:cumulative-weight}
The \emph{cumulative injection weight} of a path
$\mathcal{P} = (\beta_0,\ldots,\beta_n)$ is:
\[
  \CIW{\mathcal{P}} = \sum_{i=0}^{n} \IW{\beta_i}.
\]
\end{definition}

\begin{proposition}
\label{prop:minimal-injection-dp}
Let $\beta$ satisfy the Minimal Injection Principle, and let
$\varphi \in \Inj{\beta}$ be such that there exists an
alternative statement $\psi$ with
$\CC{\DC{\beta}} \nvdash \psi$,
$\varphi \nvdash \psi$, and $\psi \nvdash \varphi$.
Then $\beta$ is associated with a genuine choice requiring a
Decision Block.
\end{proposition}

\begin{proof}
Since $\psi$ is not derivable from prior content, injecting
$\psi$ in place of $\varphi$ represents a distinct foundational
commitment. The path $\mathcal{P}_\psi$ injecting $\psi$ and the
path $\mathcal{P}_\varphi$ injecting $\varphi$ are both locally
consistent by the Minimal Injection Principle. Since
$\varphi \nvdash \psi$ and $\psi \nvdash \varphi$, the two paths
represent genuinely different foundational commitments. By
Definition~\ref{def:genuine-choice}, this is a genuine choice.
By Axiom~\ref{ax:documentation}, it must be documented.

In the case where $\mathcal{P}_\psi \vdash \neg\varphi$, the
Decision Point is non-trivial in the sense of
Definition~\ref{def:nontrivial}. In the general case where
$\mathcal{P}_\psi \nvdash \neg\varphi$, the Decision Point is
still a genuine choice --- two locally consistent paths diverge
with different foundational commitments --- and must be
documented, even if it does not meet the stronger condition of
Definition~\ref{def:nontrivial}.
\end{proof}

\begin{remark}[Minimal Injection Applied to Logical Commitments]
\label{rem:logical-injection}
The Minimal Injection Principle applies not only to mathematical
content but to logical commitments. A logically minimal
development defaults to intuitionistic logic and introduces
classical commitments only when necessary, since intuitionistic
logic is strictly weaker than classical logic. Per-injection
logic declarations --- specifying the inference system under
which each injection is introduced --- are the logical analogue
of the per-injection necessity declarations required by
Axiom~\ref{ax:minimal}.
\end{remark}

\subsection{Architectural Transparency}
\label{subsec:transparency}

\begin{definition}[Dependency Legibility]
\label{def:dep-legible}
A path $\mathcal{P}$ is \emph{dependency-legible} if, for every
Block $\beta \in \mathcal{P}$, the set $\mathrm{Dep}(\beta)$ is
explicitly declared, and $\DC{\beta}$ is reconstructible from
these declarations without inference or supplementation.
\end{definition}

\begin{definition}[Decision Legibility]
\label{def:dec-legible}
A path $\mathcal{P}$ is \emph{decision-legible} if every
Decision Point is documented in a Decision Block satisfying
Axiom~\ref{ax:documentation}.
\end{definition}

\begin{remark}[Transparency of Decision Point Basis]
\label{rem:dec-basis}
Decision Legibility requires not only that Decision Points be
documented but that the basis of their identification ---
objective or interpretive, in the sense of
Definition~\ref{def:genuine-choice} --- be declared in the
Rationale component. A development in which Decision Points are
recorded but their basis of identification is suppressed presents
interpretive judgments as if they were formally determined.
\end{remark}

\begin{definition}[Injection Legibility]
\label{def:inj-legible}
A path $\mathcal{P}$ is \emph{injection-legible} if, for every
Block $\beta$, the injection set $\Inj{\beta}$ is explicitly
identified, every element is individually declared as an
injection, and the path satisfies Axiom~\ref{ax:minimal}.
\end{definition}

\begin{definition}[Branch Legibility]
\label{def:branch-legible}
A path $\mathcal{P}$ is \emph{branch-legible} if the Branch Tree
$\mathcal{T}$ is fully documented: every Branch is identified by
its originating Decision Point, its founding selection, and its
branching prefix; and the local consistency status of every Branch
is declared.
\end{definition}

\begin{definition}[Architectural Transparency]
\label{def:transparency}
A path $\mathcal{P}$ is \emph{architecturally transparent} if
and only if it is dependency-legible, decision-legible,
injection-legible, and branch-legible.
\end{definition}

\begin{remark}
Architectural Transparency is a conjunctive property: satisfying
three of the four conditions without the fourth yields a
development that is only partially auditable.
\end{remark}

\begin{proposition}
\label{prop:prefix-transparent}
If $\mathcal{P}$ is architecturally transparent, then every
initial segment $(\beta_0, \ldots, \beta_k)$ for $k \leq n$ is
also architecturally transparent.
\end{proposition}

\begin{proof}
Each of the four legibility conditions applies to individual
Blocks and their declared structural components. Dependency and
Injection Legibility hold for initial segments trivially. Decision
Legibility holds since every Decision Block in the initial segment
satisfies Axiom~\ref{ax:documentation}. Branch Legibility holds
since the Branch Tree of the initial segment is a sub-tree of the
full Branch Tree, and the documentation requirement for the
initial segment is satisfied as a restriction of the full
documentation.
\end{proof}

\begin{remark}[Architectural Transparency and Institution Theory]
\label{rem:institution}
Institution theory~\cite{GoguenBurstall1992} achieves a form of
logical neutrality at the mathematical level by abstracting over
the relationship between sentences and models across different
logical systems. \BOMA{} achieves neutrality by requiring that
the object logic be declared as a component of the Atomic Block.
These are different architectural strategies for the same goal.
Institution theory characterizes what is logically common across
systems; \BOMA{} characterizes what is architecturally required
of any system, regardless of its logical character. The two
frameworks are complementary.
\end{remark}

\begin{remark}[Consolidated Statement of the Architecture]
\label{rem:consolidated}
An architecturally transparent foundational development is a
Branch Tree $\mathcal{T}$ whose paths are sequences of Blocks
satisfying: explicit Block declarations (\S\ref{subsec:blocks});
acyclic dependency partial order (\S\ref{subsec:dependencies});
fully documented Decision Points (\S\ref{subsec:decision});
independently developing Branches (\S\ref{subsec:branches});
local consistency as the sole validity criterion
(\S\ref{subsec:local}); Minimal Injection for every Block
(\S\ref{subsec:injection}); and full documentation available to
any reader without reconstruction (\S\ref{subsec:transparency}).
\end{remark}

% ============================================================
\section{Branching and the Space of Foundations}
\label{sec:branching}
% ============================================================

\subsection{The Space of Foundations}
\label{subsec:foundational-space}

\begin{definition}[Foundational Space]
\label{def:foundational-space}
Let $\beta_0$ be a fixed Atomic Block. The \emph{Foundational
Space} relative to $\beta_0$, written $\FS{\beta_0}$, is the
collection of all locally consistent architectural paths whose
first element is $\beta_0$:
\[
  \FS{\beta_0}
  = \{\,\mathcal{P} \mid \mathcal{P}
    \text{ is a locally consistent architectural path with }
    \mathcal{P}_0 = \beta_0\,\}.
\]
\end{definition}

\begin{remark}
The Foundational Space is not merely a catalogue of existing
foundational developments~\cite{BenacerrafPutnam1983,vanHeijenoort1967}.
It is a \emph{possibility space}: the collection of all locally
consistent architectural paths, the vast majority of which have
not been traversed by any mathematical development.
\end{remark}

\begin{remark}[The Foundational Space as Possibility Space]
\label{rem:possibility}
Unknown mathematics is not, in this picture, mathematics that
exists somewhere waiting to be found, nor mathematics that does
not exist until invented. It is mathematics that \emph{could be}
--- paths in $\FS{\beta_0}$ whose Blocks are determinable, whose
Decision Points are documentable, and whose local consistency is
verifiable, but whose content has not yet been produced by any
development.

Every non-trivial Decision Point in an existing development is a
bifurcation: the path taken and at least one path not taken.
Every unchosen option at every documented Decision Point is an
entrance to unknown mathematics. \BOMA{} does not merely organize
what is known. It makes the structure of what is not yet known
visible, navigable, and methodologically tractable.
\end{remark}

\begin{remark}[Relation to the Set-Theoretic Multiverse]
\label{rem:multiverse}
The Foundational Space $\FS{\beta_0}$ bears a structural
resemblance to the set-theoretic multiverse proposed by
Hamkins~\cite{Hamkins2012}, and the relationship between the two
frameworks merits explicit discussion.

Hamkins argues that the independence results of set theory
motivate a pluralist ontological stance: there is no single
``true'' set-theoretic universe, but rather a multiverse of
equally legitimate universes, each with its own truths.

The Foundational Space $\FS{\beta_0}$ differs from Hamkins'
multiverse in five fundamental respects.

\textit{Scope.} Hamkins' multiverse is exclusively set-theoretic.
The Foundational Space accommodates paths of any foundational
character --- set-theoretic, type-theoretic, categorical, or
otherwise. The set-theoretic paths in $\FS{\beta_0}$ constitute a
specific cluster analogous to Hamkins' multiverse, but a proper
substructure of the full Foundational Space.

\textit{Level of operation.} Hamkins' multiverse operates at the
level of mathematical models, with relationships established by
forcing and inner models. The Foundational Space operates at the
level of architectural construction, with relationships
established by common ancestry, divergence points, and Decision
Point reselections.

\textit{Ontological commitment.} Hamkins' multiverse carries an
explicit ontological commitment: the plurality of universes is a
claim about what genuinely exists. The Foundational Space carries
no such commitment.

\textit{Documentation.} The universes of Hamkins' multiverse
exist independently of any documentation of the decisions that
produced them. In contrast, a path in $\FS{\beta_0}$ is an
architecturally transparent development with explicitly
documented Decision Points and injection structure.

\textit{Mechanism of plurality.} In Hamkins' multiverse, plurality
arises from model-theoretic operations within set theory itself.
In the Foundational Space, plurality arises from architectural
decisions that precede mathematical development.

In summary: Hamkins' multiverse can be understood as a specific,
densely explored region of the Foundational Space --- the cluster
of set-theoretic paths diverging at Decision Points internal to
ZFC and its extensions --- studied through the lens of
model-theoretic semantics rather than architectural construction.
\end{remark}

\subsection{Structure of the Foundational Space}
\label{subsec:fs-structure}

\begin{definition}[Common Ancestry]
\label{def:common-ancestry}
The \emph{common ancestry} of $\mathcal{P}$ and $\mathcal{Q}$,
written $\CA{\mathcal{P}}{\mathcal{Q}}$, is the longest initial
segment shared by both paths.
\end{definition}

\begin{definition}[Divergence Point]
\label{def:divergence}
The \emph{divergence point} of $\mathcal{P}$ and $\mathcal{Q}$
is the Decision Point $\mathcal{D}$ immediately following their
common ancestry.
\end{definition}

\begin{proposition}[Architectural Ultra-Pseudometric]
\label{prop:metric}
Define $d : \FS{\beta_0} \times \FS{\beta_0} \to \mathbb{R}_{\geq 0}$ by:
\[
  d(\mathcal{P}, \mathcal{Q})
  = \frac{1}{|\CA{\mathcal{P}}{\mathcal{Q}}| + 1},
\]
where $|\CA{\mathcal{P}}{\mathcal{Q}}|$ is taken as $\infty$ when
the common ancestry is infinite, with convention
$\frac{1}{\infty+1}=0$. Then $d$ satisfies non-negativity,
symmetry, and the ultrametric inequality, constituting an
\emph{ultra-pseudometric} on $\FS{\beta_0}$.
\end{proposition}

\begin{proof}
Non-negativity and symmetry are immediate. For the ultrametric
inequality: let $j = |\CA{\mathcal{P}}{\mathcal{Q}}|$ and
$k = |\CA{\mathcal{Q}}{\mathcal{R}}|$, and set $m = \min(j,k)$.
Since $\mathcal{P}$ and $\mathcal{Q}$ agree on their first $j$
Blocks, and $\mathcal{Q}$ and $\mathcal{R}$ agree on their first
$k$ Blocks, both $\mathcal{P}$ and $\mathcal{R}$ agree with
$\mathcal{Q}$ on their first $m$ Blocks, hence agree with each
other on those $m$ Blocks. Therefore
$|\CA{\mathcal{P}}{\mathcal{R}}| \geq m = \min(j,k)$, giving:
\[
  d(\mathcal{P},\mathcal{R})
  \leq \frac{1}{\min(j,k)+1}
  = \max\!\left(\frac{1}{j+1},\frac{1}{k+1}\right)
  = \max(d(\mathcal{P},\mathcal{Q}),\,d(\mathcal{Q},\mathcal{R})).
\qedhere
\]
\end{proof}

\begin{remark}[Packaging Dependence and Its Interpretation]
\label{rem:packaging}
The ultra-pseudometric $d$ depends on $|\CA{\mathcal{P}}{\mathcal{Q}}|$
--- the Block count of the common ancestry --- rather than on any
intrinsic measure of logical content. The same logical content
packaged into more Blocks yields a smaller distance.

This dependence is a feature, not a defect, provided its
interpretation is correctly understood. The function $d$ measures
\emph{architectural proximity} --- the degree to which two paths
proceeded in documented agreement --- rather than \emph{logical
proximity}. Architectural proximity is the appropriate measure in
the context of \BOMA{}, which is concerned with the process of
foundational construction and not only with the resulting content.
\end{remark}

\begin{remark}[Pseudometric Character]
\label{rem:pseudometric}
The function $d$ fails to be a genuine metric when infinite paths
are admitted, since two paths sharing an infinite common ancestry
receive distance zero without being identical. This correctly
reflects architectural indistinguishability at any finite depth.
For all practical purposes --- finite foundational developments
--- the restriction to finite-length paths gives a genuine
ultrametric.
\end{remark}

\begin{remark}[Canonical Normalization]
\label{rem:normalization}
When the primary purpose of computing $d$ is to compare the
logical depth of shared commitment, paths may be normalized to a
common level of granularity before computing the Block count. A
\emph{canonical normalization} re-packages content into Blocks of
uniform logical weight (one injection per Block). The normalized
distance $d^*$ is invariant under re-packaging and measures shared
logical commitment. The relationship between $d$ and $d^*$ is
itself architecturally informative.

Canonical normalization in \S\ref{subsec:comparison-multi} provides a
basis for comparison invariant under packaging decisions.
\end{remark}

\begin{definition}[Ancestry Depth]
\label{def:ancestry-depth}
The \emph{ancestry depth} of $\mathcal{P}$ and $\mathcal{Q}$
is $\delta(\mathcal{P},\mathcal{Q}) =
|\CA{\mathcal{P}}{\mathcal{Q}}|$.
\end{definition}

\subsection{Navigation in the Foundational Space}
\label{subsec:navigation}

\begin{definition}[Reachability]
\label{def:reachability}
A path $\mathcal{Q}$ is \emph{reachable} from $\mathcal{P}$ if
there exists a finite sequence of Decision Point reselections
connecting them. Since the Branch Tree is connected, every pair
of paths is mutually reachable.
\end{definition}

\subsection{Contradictions as Structural Information}
\label{subsec:contradictions}

\begin{definition}[Inter-path Contradiction]
\label{def:inter-contradiction}
An \emph{inter-path contradiction} relative to $\mathcal{D}$ is
a formula $\varphi$ such that
$\mathcal{P} \vdash \varphi$ and $\mathcal{Q} \vdash \neg\varphi$
for branches $\mathcal{P}$, $\mathcal{Q}$ diverging at
$\mathcal{D}$.
\end{definition}

\begin{theorem}[Contradiction Localization]
\label{thm:localization}
Let $\mathcal{P}$ and $\mathcal{Q}$ be locally consistent paths
with $\mathcal{P} \vdash \varphi$ and
$\mathcal{Q} \vdash \neg\varphi$. Then:
\begin{enumerate}
  \item Neither $\mathcal{P}$ nor $\mathcal{Q}$ is internally
        inconsistent.
  \item The inter-path contradiction is fully localized to the
        divergence point $\mathcal{D}$.
  \item The inter-path contradiction certifies that $\mathcal{D}$
        is non-trivial.
\end{enumerate}
\end{theorem}

\begin{proof}
(1) Both paths are locally consistent by hypothesis.

(2) By Proposition~\ref{prop:inheritance}, the cumulative content
of $\CA{\mathcal{P}}{\mathcal{Q}}$ is identical in both paths.

\textit{That $\varphi$ is not derivable from common ancestry:}
Suppose $\CA{\mathcal{P}}{\mathcal{Q}} \vdash \varphi$. Then both
paths derive $\varphi$, but $\mathcal{Q} \vdash \neg\varphi$,
so $\mathcal{Q}$ derives both $\varphi$ and $\neg\varphi$,
contradicting local consistency of $\mathcal{Q}$.

\textit{That $\neg\varphi$ is not derivable from common ancestry:}
Suppose $\CA{\mathcal{P}}{\mathcal{Q}} \vdash \neg\varphi$. Then
both paths derive $\neg\varphi$, but $\mathcal{P} \vdash \varphi$,
so $\mathcal{P}$ derives both $\varphi$ and $\neg\varphi$,
contradicting local consistency of $\mathcal{P}$.

Hence the derivations must depend on post-Decision content
originating at $\mathcal{D}$.

(3) Immediate from Definition~\ref{def:nontrivial}.
\end{proof}

\begin{definition}[Separation Formula and Semantic Signature]
\label{def:separation}
A \emph{separation formula} for $\mathcal{D}$ is a formula
$\varphi \in \Sep{\mathcal{D}}$ such that
$\mathcal{P} \vdash \varphi$ and $\mathcal{Q} \vdash \neg\varphi$
for two locally consistent branches diverging at $\mathcal{D}$.
The set $\Sep{\mathcal{D}}$ is the \emph{semantic signature} of
$\mathcal{D}$.
\end{definition}

\subsection{Mathematics as a Traceable Knowledge System}
\label{subsec:traceable}

\begin{definition}[Traceable Knowledge System]
\label{def:traceable}
A foundational development is a \emph{traceable knowledge system}
if every statement in its deductive closure can be assigned a
\emph{provenance} --- a complete record of the Blocks and Decision
Point selections from which it ultimately derives --- and this
provenance is structurally determined from the architectural record.
\end{definition}

\begin{proposition}
\label{prop:traceable}
Every architecturally transparent path in the sense of
Definition~\ref{def:transparency} is a traceable knowledge system.
\end{proposition}

\begin{proof}
Architectural Transparency requires Dependency Legibility,
Injection Legibility, and Decision Legibility. These conditions
ensure that the dependency closure of every Block, the injection
set of every Block, and the selection at every Decision Point are
explicitly declared. A derivation of $\varphi$ in $\mathcal{P}$
proceeds through a finite sequence of Blocks; its provenance is
the subgraph of the dependency structure supporting this
derivation, together with the injections and Decision Point
selections on which it depends. This provenance is structurally
determined --- a definite object fixed by the architectural record.
\end{proof}

\begin{remark}[On Computability and Structural Determination]
\label{rem:computability-provenance}
The claim is that provenance is \emph{structurally determined},
not that it is \emph{mechanically extractable} in all cases.
Extraction depends on the decidability properties of the
underlying logic. In cases where the logic has a decidable
consequence relation, provenance extraction is mechanically
feasible. In the general first-order case, it requires
non-mechanical means.
\end{remark}

% ============================================================
\section{Comparison as an Epistemological Tool}
\label{sec:comparison}
% ============================================================

\subsection{What Comparison Is Not}
\label{subsec:not-comparison}

\begin{definition}[Adjudicative Comparison]
\label{def:adjudicative}
An \emph{adjudicative comparison} of two paths produces a
judgment of the form: $\mathcal{P}$ is preferable to
$\mathcal{Q}$, where preferability is defined in terms of
validity, correctness, or foundational superiority.
\end{definition}

\begin{principle}[Prohibition on Adjudicative Comparison]
\label{prin:no-adjudication}
\BOMA{} provides no mechanism for adjudicative comparison. No
structural feature --- not the injection weight, not the depth of
common ancestry, not the cardinality of the semantic signature
--- constitutes a criterion of foundational superiority between
locally consistent paths.
\end{principle}

\subsection{What Comparison Is}
\label{subsec:what-comparison}

\begin{definition}[Structural Comparison]
\label{def:structural-comparison}
A \emph{structural comparison} of $\mathcal{P}$ and $\mathcal{Q}$
produces: the common ancestry; the divergence point; the
selections made at the divergence point in each path; an
approximation of the semantic signature; the injection weights
of post-Decision Blocks; and the local consistency status of
each path. No output includes a judgment of preferability.
\end{definition}

\begin{definition}[Epistemological Yield]
\label{def:yield}
The \emph{epistemological yield} of a structural comparison is
the knowledge produced about the architecture of the Foundational
Space --- specifically about the nature of the divergence point,
the mathematical consequences of the selection, and the
relationship between the paths within $\FS{\beta_0}$.
\end{definition}

\subsection{Structural Ancestry and the Locus of Difference}
\label{subsec:ancestry}

\begin{proposition}[Difference Localization]
\label{prop:diff-localize}
Every mathematical difference between $\mathcal{P}$ and
$\mathcal{Q}$ is attributable to post-Decision content and
localized to the divergence point $\mathcal{D}$.
\end{proposition}

\begin{proof}
By Proposition~\ref{prop:inheritance}, cumulative content of
$\CA{\mathcal{P}}{\mathcal{Q}}$ is identical in both paths.
Any statement derivable from common ancestry is derivable in
both. Any disagreement cannot derive from common ancestry and
must depend on post-Decision content originating at $\mathcal{D}$.
\end{proof}

\begin{definition}[Injection Divergence]
\label{def:inj-divergence}
The \emph{injection divergence} of $\mathcal{P}$ and $\mathcal{Q}$
with divergence point $\mathcal{D}$ is:
\[
  \Delta_{\mathrm{Inj}}(\mathcal{P},\mathcal{Q})
  = \Bigl(
      \bigcup_{\beta \in \mathcal{P}\setminus\CA{\mathcal{P}}{\mathcal{Q}}}
        \Inj{\beta}
    \Bigr)
    \;\triangle\;
    \Bigl(
      \bigcup_{\beta \in \mathcal{Q}\setminus\CA{\mathcal{P}}{\mathcal{Q}}}
        \Inj{\beta}
    \Bigr).
\]
\end{definition}

\subsection{What We Learn from Contradictory Branches}
\label{subsec:learn}

\begin{theorem}[Epistemic Content of Contradiction]
\label{thm:epistemic}
Let $\mathcal{P}$ and $\mathcal{Q}$ be locally consistent paths
with divergence point $\mathcal{D}$ and separation formula
$\varphi \in \Sep{\mathcal{D}}$. The structural comparison of
$\mathcal{P}$ and $\mathcal{Q}$ with respect to $\varphi$ yields:
\begin{enumerate}
  \item \textbf{Consistency confirmation}: both paths are locally
        consistent;
  \item \textbf{Decision Point certification}: $\mathcal{D}$ is
        non-trivial;
  \item \textbf{Assumption identification}: $\Delta_{\mathrm{Inj}}
        (\mathcal{P},\mathcal{Q})$ contains the assumptions
        responsible for the different derivations;
  \item \textbf{World demarcation}: $\varphi$ is a statement
        whose truth value is determined by the selection at
        $\mathcal{D}$;
  \item \textbf{Irreducibility confirmation}: the disagreement
        is not resolvable within \BOMA{}.
\end{enumerate}
\end{theorem}

\begin{proof}
Items (1) and (2) follow from Theorem~\ref{thm:localization}.
Item (3) follows from Definition~\ref{def:inj-divergence} and
Axiom~\ref{ax:minimal}: since every injection is individually
declared and irredundant, assumptions responsible for any derived
statement are identifiable from injection sets.
Item (4) follows from Definition~\ref{def:separation}.
Item (5) follows from Principle~\ref{prin:no-adjudication}.
\end{proof}

\subsection{Comparison Across Multiple Paths}
\label{subsec:comparison-multi}

\begin{definition}[Comparative Profile]
\label{def:comparative-profile}
Let $\{\mathcal{P}_1,\ldots,\mathcal{P}_m\}$ be locally
consistent paths in $\FS{\beta_0}$. Their \emph{comparative
profile} consists of: for each pair, the common ancestry,
divergence point, and ancestry depth; the pairwise distance
matrix $M_{ij} = d(\mathcal{P}_i,\mathcal{P}_j)$; for each
Decision Point, the partition of paths by their selections;
and for each separation formula, the partition of paths into
those deriving $\varphi$ and those deriving $\neg\varphi$.

Canonical normalization of paths --- as defined in
Remark~\ref{rem:normalization} --- provides a basis for
comparison that is invariant under packaging decisions.
\end{definition}

\begin{definition}[Foundational Cluster]
\label{def:cluster}
A subset $\mathcal{C} \subseteq \{\mathcal{P}_1,\ldots,\mathcal{P}_m\}$
is a \emph{foundational cluster at depth $k$} if
$\delta(\mathcal{P}_i,\mathcal{P}_j) \geq k$ for all
$\mathcal{P}_i,\mathcal{P}_j \in \mathcal{C}$.
\end{definition}

\begin{remark}
The ancestry depth of Definition~\ref{def:ancestry-depth} is
related to the branch depth of Definition~\ref{def:branch-depth}:
two paths diverging at a Branch of depth $d$ have ancestry depth
equal to the number of Blocks in the branching prefix at that
depth. Branches at greater depth share longer common ancestry.
\end{remark}

\subsection{Comparison as a Research Methodology}
\label{subsec:methodology}

\begin{definition}[Architectural Inquiry]
\label{def:architectural-inquiry}
An \emph{architectural inquiry} is a research program that
applies structural comparison to a collection of paths with
the goal of producing their comparative profile, identifying
non-trivial Decision Points, characterizing their semantic
signatures, and mapping the foundational clusters that emerge.
\end{definition}

% ============================================================
\section{Relation to Existing Foundations}
\label{sec:existing}
% ============================================================

\begin{remark}[Architectural Reconstruction]
\label{rem:reconstruction}
The analyses of existing foundational systems in this section
employ \emph{architectural reconstruction}: the retrospective
reading of an existing foundational development through the lens
of \BOMA{}, with the aim of identifying structures corresponding
to Blocks, dependency relations, Decision Points, and injection
sets.

Architectural reconstruction is an interpretive practice, not a
mechanical extraction. The Decision Points identified below are
not historical events in the sense of discrete, localized
occurrences. They are logical structures --- the existence of
mutually consistent alternatives from which a selection was
eventually made --- not specific moments at which specific persons
made specific choices.

The primary purpose of \BOMA{} is prospective: it is designed for
future foundational developments conducted within the
architectural paradigm. The application to existing foundations
is a secondary, illustrative purpose.
\end{remark}

\subsection{Zermelo-Fraenkel Set Theory with Choice}
\label{subsec:zfc}

\begin{observation}[ZFC as an Architectural Path]
\label{obs:zfc-path}
ZFC~\cite{Jech2003,Zermelo1908} corresponds to a locally
consistent architectural path $\mathcal{P}_{\mathrm{ZFC}}$ in
$\FS{\beta_0^{\mathrm{ZFC}}}$, where the Atomic Block declares:
the membership relation $\in$ as primitive; classical first-order
logic; the universe of discourse consisting entirely of sets.
Every axiom of ZFC corresponds to an injecting Block.
\end{observation}

\begin{observation}[Decision-Point-Structured Episodes in ZFC]
\label{obs:zfc-episodes}
The historical development of ZFC contains several episodes with
the logical structure of Decision Points.

\textit{The Choice Episode.} Between 1904 and the mid-twentieth
century~\cite{BaireEtAl1905,Moore1982,Zermelo1904}, the
mathematical community engaged in sustained deliberation about
the Axiom of Choice. The episode involved explicit recognition
of alternatives --- Choice, Dependent Choice, Countable Choice,
no choice principle --- and explicit argument about their
consequences. The architectural framework reads this as a
Decision Point with semantic signature including the
Well-Ordering Theorem~\cite{Zermelo1904}, Zorn's
Lemma~\cite{Jech2003}, and the Hahn-Banach
Theorem~\cite{Kunen1980}.

\textit{The Foundation Episode.} The Axiom of Foundation was
introduced primarily as a technical device excluding pathological
constructions. The alternative --- non-well-founded sets ---
has been developed by Aczel and others~\cite{Aczel1988,AczelRathjen2001}.
The reconstruction of this as a Decision Point is more tentative:
the ``choice'' was partly a decision by default.

\textit{The Classical Logic Episode.} The commitment to classical
first-order logic was made against a background in which
intuitionistic alternatives were known --- Brouwer's
intuitionism~\cite{Brouwer1907} and Heyting's
formalization~\cite{Heyting1930,Heyting1956} were contemporaneous.
\end{observation}

\begin{remark}
The identification of Decision-Point-structured episodes in ZFC
is an architectural reconstruction in the sense of
Remark~\ref{rem:reconstruction}. It is an illustration of the
framework's concepts, not a historical claim that ZFC consciously
proceeded through documented Decision Points.
\end{remark}

\begin{remark}[\BOMA{} and the Set-Theoretic Multiverse]
\label{rem:multiverse-zfc}
Within the set-theoretic region of the Foundational Space, the
Decision Points corresponding to large cardinal
axioms~\cite{Jech2003,Kunen1980}, forcing-related
extensions~\cite{Cohen1963,Cohen1964}, and variants of the Axiom
of Choice~\cite{Moore1982} generate a cluster of paths whose
structure is analogous to Hamkins' set-theoretic
multiverse~\cite{Hamkins2012}.

The architectural framework adds two elements that Hamkins'
multiverse does not contain. First, the architectural record: the
Decision Point at which the large cardinal axiom was injected,
the alternatives, and the rationale are documented within the
path. Second, the architectural distance: the ultra-pseudometric
of Proposition~\ref{prop:metric} provides a measure of how far
any two set-theoretic paths are from their common ancestry.
\end{remark}

\subsection{Peano Arithmetic and Its Extensions}
\label{subsec:pa}

\begin{observation}[Peano Arithmetic as an Architectural Path]
\label{obs:pa-path}
Peano Arithmetic~\cite{Peano1889} corresponds to a locally
consistent path $\mathcal{P}_{\mathrm{PA}}$ whose Atomic Block
declares: a distinguished element $0$, a unary successor function
$S$, and equality; classical first-order logic; the natural
numbers as domain.
\end{observation}

\begin{observation}[Decision Points in the Peano Development]
\label{obs:pa-episodes}
The most consequential episode is the \emph{Induction Strength
Episode}. Options include: first-order induction~\cite{Peano1889};
second-order induction, yielding $Z_2$~\cite{Simpson2009}; and
intermediate systems studied in reverse
mathematics~\cite{Friedman1975,Simpson2009}.
\end{observation}

\begin{remark}
The program of reverse mathematics, developed by Friedman~\cite{Friedman1975}
and extended by Simpson~\cite{Simpson2009}, is, from the
architectural perspective, a systematic investigation of a
specific region of the Foundational Space. Its central question
--- which axioms are necessary and sufficient for specific
mathematical theorems? --- is an architectural question.
\end{remark}

\subsection{Dependent Type Theory}
\label{subsec:type-theory}

\begin{observation}[Type Theory as an Architectural Path]
\label{obs:mltt-path}
Martin-L\"{o}f Type Theory~\cite{MartinLof1975,MartinLof1984}
corresponds to a locally consistent path $\mathcal{P}_{\mathrm{MLTT}}$
whose Atomic Block declares: types and terms as primitive, with
judgment forms $A\,\mathrm{type}$, $a : A$, $A = B\,\mathrm{type}$,
and $a = b : A$; the propositions-as-types interpretation.
\end{observation}

\begin{observation}[Decision Points Within Type Theory]
\label{obs:tt-episodes}
The \emph{Extensionality Episode} concerns whether the identity
type $\mathrm{Id}_A(a,b)$~\cite{MartinLof1984} is required to
satisfy uniqueness of identity proofs~\cite{MartinLof1975}. The
selection of the non-uniqueness option yields intensional type
theory, which is the foundation on which HoTT is built.

The \emph{Universe Hierarchy Episode} concerns whether the type
theory includes a cumulative hierarchy $\mathcal{U}_0 :
\mathcal{U}_1 : \ldots$ Different selections produce different
type theories with different expressive powers.
\end{observation}

\subsection{Homotopy Type Theory}
\label{subsec:hott}

\begin{observation}[HoTT as an Architectural Branch]
\label{obs:hott-branch}
HoTT~\cite{HoTTBook2013,Voevodsky2010} corresponds to a locally
consistent path branching from $\mathcal{P}_{\mathrm{MLTT}}$ at
the \emph{Univalence Episode}. The Univalence Axiom is not
derivable within intensional MLTT and therefore represents a new
foundational commitment. Separation formulas for this Decision
Point include statements provable in HoTT but not in MLTT.
\end{observation}

\subsection{Categorical Foundations}
\label{subsec:categorical}

\begin{observation}[Categorical Foundations as Architectural Paths]
\label{obs:etcs-path}
ETCS, developed by Lawvere~\cite{Lawvere1964}, corresponds to a
locally consistent path $\mathcal{P}_{\mathrm{ETCS}}$ whose Atomic
Block declares: objects, morphisms, composition, and identity as
primitive~\cite{MacLane1971,Awodey2010}; sets understood as objects
of a specific category rather than extensionally determined
collections.
\end{observation}

\begin{remark}
It is known that ETCS and a bounded version of ZFC are mutually
interpretable~\cite{Lawvere1964,McLarty1991}. Within the
architectural framework, this mutual interpretability corresponds
to a structural relationship between the two paths --- a map
between their deductive closures preserving logical consequence in
both directions. This does not collapse the architectural
distinction between the two paths.
\end{remark}

\subsection{The Architectural Status of Existing Foundations}
\label{subsec:arch-status}

\begin{proposition}[Existing Foundations as Architecturally Opaque Paths]
\label{prop:opaque}
Every existing foundational system corresponds to a locally
consistent path in an appropriate Foundational Space. Every such
path is valid in the sense of Definition~\ref{def:validity}. No
existing foundational system is rendered invalid, superseded, or
diminished by \BOMA{}.
\end{proposition}

\begin{proof}
Local consistency of each system follows from its consistency.
The remaining claims follow from the definitions: every locally
consistent path is a point in the Foundational Space, and \BOMA{}
places no restriction on content beyond local consistency and the
general structural requirements of Definition~\ref{def:path}.
\end{proof}

\begin{proposition}[Existing Foundations as Special Cases]
\label{prop:special-cases}
Every existing foundational system, viewed as an architectural
path, satisfies the structural requirements of
Definition~\ref{def:path} but does not satisfy Architectural
Transparency as defined in Definition~\ref{def:transparency}.
Specifically, existing foundations are decision-opaque: their
Decision Points are not documented within the formal presentation.
\end{proposition}

\begin{remark}
Architectural opacity is not a defect; it is a property that
follows from the design goals of the existing paradigm. The
application of \BOMA{} to existing foundations is an instance of
architectural reconstruction in the sense of
Remark~\ref{rem:reconstruction}: a retrospective interpretive
reading. Existing foundations remain mathematically valid,
internally consistent, and enormously productive.
\end{remark}

\subsection{Compatibility and Non-Competition}
\label{subsec:compatibility}

\begin{principle}[Non-Competition]
\label{prin:noncompetition}
\BOMA{} does not compete with any existing foundational system.
It does not claim that ZFC is incorrect, that Peano Arithmetic is
inadequate, that type theories are misguided, or that categorical
foundations are inferior to any alternative.
\end{principle}

\begin{principle}[Additive Extension]
\label{prin:additive}
\BOMA{} is an \emph{additive extension} of the existing
foundational paradigm: it adds the structural dimension of
documented Decision Points, explicit injection structure, and
maintained Branch Trees, without removing or altering any content
that existing presentations provide.
\end{principle}

\subsection{The Architecture as a Meta-Framework}
\label{subsec:meta}

\begin{definition}[Meta-Framework]
\label{def:meta}
A \emph{meta-framework} for foundational mathematics is a
structure within which individual foundational systems can be
situated, compared, and related to one another, without itself
constituting a foundational system.
\end{definition}

\begin{proposition}
\label{prop:meta}
\BOMA{} is a meta-framework in the sense of
Definition~\ref{def:meta} with respect to the \emph{object logic}
of individual paths: it imposes no constraint on which logic any
particular path declares. At the level of its own formulation,
\BOMA{} employs classical first-order logic as its framework
logic --- a commitment that is explicit and acknowledged as a
meta-level Decision Point in the sense of
Remark~\ref{rem:three-levels}.
\end{proposition}

\begin{proof}
The framework makes no specific mathematical commitments at the
object level. The additional qualification --- that logical
neutrality holds at the object level but not at the framework or
meta-logic levels --- follows from the explicit use of classical
reasoning in the framework's own definitions and proofs, as
established in Remark~\ref{rem:three-levels}.
\end{proof}

\begin{remark}[The Meta-Level Decision Point]
\label{rem:meta-dp}
The choice of classical logic as the framework logic is, by
\BOMA{}'s own standards, a Decision Point. The present paper
selects classical logic on grounds of familiarity, expressive
adequacy, and consistency with the conventions of the mathematical
logic community. The unchosen alternative --- an intuitionistic
formulation of the same architectural framework --- is preserved
in the record by this acknowledgment.
\end{remark}

% ============================================================
\section{Applications and Research Program}
\label{sec:applications}
% ============================================================

\subsection{Formalization and Proof Assistants}
\label{subsec:proof-assistants}

\begin{observation}[The Current State of Proof Assistant Foundations]
\label{obs:pa-current}
Existing proof assistants are grounded in specific foundational
systems: Coq~\cite{Coq2023} and Agda~\cite{Norell2007} in
variants of MLTT; Lean~\cite{deMouraEtAl2015} in a dependent
type theory; Isabelle~\cite{NipkowEtAl2002} in Higher-Order
Logic; Metamath~\cite{Megill1997} in a metalogical framework.
Each of these is, in the architectural sense, an architecturally
opaque path.
\end{observation}

\begin{application}[Architecturally Transparent Proof Assistants]
\label{app:transparent-pa}
\BOMA{} suggests a new design requirement for proof assistants:
not merely verification of formal derivations, but documentation
of the architectural choices made in the underlying foundational
development, maintained as a first-class structural artifact.
\end{application}

\begin{application}[Dependency Analysis and Axiom Tracking]
\label{app:axiom-tracking}
Given an architecturally transparent development, the injection
set of every theorem --- the collection of assumptions on which
it ultimately depends --- is structurally determined by the
architectural record. In cases where the underlying logic supports
effective proof search, this injection set is mechanically
extractable. This enables \emph{axiom profiling}: for any theorem,
a precise account of which foundational commitments are necessary
for its proof.

A user of a large formal library --- such as Mathlib in
Lean~\cite{LeanCommunity2023}, or the Coq standard
library~\cite{Coq2023} --- who wishes to understand why a
specific definition was chosen over an alternative currently has
no structural source for this information within the library.
An architecturally transparent library would make this a
first-class structural artifact.
\end{application}

\begin{application}[Library Branching and Foundational Experiments]
\label{app:branching}
The Branch structure of \BOMA{} suggests a new mode of
interaction with formal libraries: the ability to branch a library
at a specific Decision Point and explore the consequences of a
different foundational choice, while preserving the original
development unchanged.
\end{application}

\subsection{Philosophy of Mathematics}
\label{subsec:philosophy}

\begin{application}[A Formal Framework for Foundational Pluralism]
\label{app:pluralism}
The Foundational Space $\FS{\beta_0}$ provides a formal framework
for foundational pluralism. Within \BOMA{}, foundational pluralism
is not a philosophical thesis to be argued for or against; it is
a structural fact.
\end{application}

\begin{application}[Independence Results as Architectural Information]
\label{app:independence}
The independence results of mathematical logic --- G\"{o}del's
incompleteness theorems~\cite{Goedel1931}, the independence of
the Continuum Hypothesis from ZFC~\cite{Cohen1963,Cohen1964,Goedel1938,Goedel1940}
--- have a precise architectural interpretation. A statement
$\varphi$ independent of a foundation $\mathcal{P}$ is a
separation formula for the Decision Point at which $\mathcal{P}$
must make a choice about $\varphi$. Independence results are
architectural certifications of non-trivial Decision Points.
\end{application}

\subsection{Mathematical Education}
\label{subsec:education}

\begin{application}[Teaching Mathematics Through Decision Points]
\label{app:education}
The architectural discipline suggests presenting mathematical
content through its Decision Points: making visible at each stage
the choices that were made and the alternatives that existed.
\end{application}

\begin{application}[Curriculum Design and Dependency Transparency]
\label{app:curriculum}
An architecturally transparent mathematical development provides
a precise dependency structure from which a well-ordered
curriculum can be read directly. The partial order on Blocks
determines logically admissible orderings of curriculum topics.
\end{application}

\subsection{Foundations Research}
\label{subsec:foundations-research}

\begin{application}[Systematic Mapping of the Foundational Space]
\label{app:mapping}
The Foundational Space $\FS{\beta_0}$ is an object of mathematical
study in its own right. Systematic mapping of specific regions
--- for example, the region containing classical and constructive
set theories --- is a research program that \BOMA{} makes coherent
and tractable.
\end{application}

\begin{application}[Classification of Decision Points]
\label{app:classification}
Decision Points vary in the depth and breadth of their
mathematical consequences. Their classification by semantic
signatures, independence relations, and architectural coupling
constitutes a theory of foundational choice.
\end{application}

\begin{application}[Minimal Foundations]
\label{app:minimal}
The Minimal Injection Principle suggests identifying minimal
foundations: paths that achieve a specified mathematical result
with the smallest possible cumulative injection weight. This
extends the reverse mathematics of Friedman~\cite{Friedman1975}
and Simpson~\cite{Simpson2009} to the full Foundational Space.
\end{application}

\begin{application}[Formal Comparison of Foundational Systems]
\label{app:formal-comparison}
The structural comparison framework of \S\ref{sec:comparison}
enables comparison at the level of architectural structure rather
than interpretability or consistency strength.

The notion of a Minimal Branch --- as defined in
Definition~\ref{def:minimal-branch} --- is particularly relevant
here: a minimal Branch isolates the effect of a single Decision
Point selection as cleanly as possible, making the structural
comparison maximally informative.
\end{application}

\subsection{Connections to Adjacent Fields}
\label{subsec:adjacent}

\begin{connection}[Software Engineering and Version Control]
\label{con:version-control}
The architectural framework is structurally analogous to
version-controlled software development~\cite{Coq2023,deMouraEtAl2015}.
Blocks correspond to commits; dependency structures correspond to
commit histories; Decision Points correspond to branch points;
Branches correspond to repository branches. The formal
infrastructure of version control systems --- directed acyclic
graphs of commits, branch management --- is a computational
realization of structures closely related to those defined in
\BOMA{}.
\end{connection}

\begin{connection}[Formal Epistemology]
\label{con:epistemology}
The notion of a traceable knowledge system (Definition~\ref{def:traceable})
connects \BOMA{} to formal epistemology. A traceable knowledge
system makes provenance a formally defined structural concept,
providing a mathematical model for one of the central concepts
of formal epistemology.
\end{connection}

\begin{connection}[Artificial Intelligence and Mathematical Reasoning]
\label{con:ai}
A system trained on architecturally transparent mathematical
content would have access to a richer structural representation:
not merely the content of mathematical developments, but the
architectural choices that shaped them and the alternatives that
were available.
\end{connection}

\subsection{The Immediate Research Agenda}
\label{subsec:agenda}

\begin{description}
  \item[Formal Verification of the Framework.]
    The encoding of \BOMA{}'s components in a proof assistant and
    the mechanically verified proof of its main propositions.
    This would produce the first architecturally transparent
    formal development --- a development of \BOMA{} itself within
    the architectural paradigm.

  \item[Reconstruction of a Classic Foundation.]
    The reconstruction of ZFC or Peano Arithmetic as an
    architecturally transparent path, producing the first
    systematic Decision Point Catalogue for a major foundational
    system.

  \item[Development of Comparison Tools.]
    Algorithms for computing common ancestries, divergence points,
    semantic signatures, and injection divergences from structural
    records.

  \item[First Branch Developments.]
    The systematic exploration of what follows from alternative
    selections at specific, architecturally characterized Decision
    Points --- beginning with the Decision Points identified in
    \S\ref{subsec:zfc} through \S\ref{subsec:categorical}.
\end{description}

% ============================================================
\section{Conclusion}
\label{sec:conclusion}
% ============================================================

\subsection{What This Paper Has Not Done}
\label{subsec:not-done}

This paper has not proposed a new mathematical foundation. It has
not introduced new axioms, defined new primitive notions, or
derived new mathematical theorems. It has not claimed that
existing foundations are deficient. It has not resolved the
philosophical questions surrounding foundational pluralism. It
has not developed any Branch in full. The architecture has been
defined; its application is the work of subsequent research.

\subsection{What This Paper Has Done}
\label{subsec:done}

We have identified an architectural gap in the existing
foundational paradigm: not a logical gap, but an organizational
one, concerning the documentation of Decision Points, the
preservation of alternatives, and the declaration of injection
structure.

We have proposed \BOMA{}, a precise architectural framework
addressing this gap through seven structural components.

We have introduced the Foundational Space as the global object
produced by the architectural framework, established
inter-path contradictions as structural information
(Theorem~\ref{thm:localization}), and characterized
architecturally transparent paths as traceable knowledge systems
(Proposition~\ref{prop:traceable}).

We have established comparison as an epistemological tool
(Theorem~\ref{thm:epistemic}) and demonstrated that \BOMA{} is a
meta-framework within which existing foundations ---
ZFC~\cite{Jech2003,Zermelo1908},
Peano Arithmetic~\cite{Peano1889},
Martin-L\"{o}f Type Theory~\cite{MartinLof1975,MartinLof1984},
HoTT~\cite{HoTTBook2013}, and categorical
foundations~\cite{Lawvere1964,MacLane1971} --- are situated as
locally consistent, architecturally opaque paths.

\subsection{The Central Claim}
\label{subsec:central}

The central claim of this paper is the following.

The organizational paradigm of existing foundational mathematics
--- linear presentation of axioms, definitions, and theorems,
with Decision Points invisible, alternatives suppressed, and
injection structure undeclared --- is one paradigm among possible
paradigms~\cite{BenacerrafPutnam1983,Shapiro2000}, and not the
only one consistent with mathematical rigor.

\BOMA{} is an alternative paradigm: one in which every
foundational development is organized as a Branch Tree of locally
consistent paths, every genuine choice is recorded at the moment
it is made, every new assumption is individually declared, and
every development is situated within a navigable space of
alternative developments.

\subsection{What Comes Next}
\label{subsec:next}

The immediate next steps are those identified in
\S\ref{subsec:agenda}: formal verification, reconstruction of a
classic foundation, development of comparison tools, and first
Branch developments.

\begin{remark}[Scope and the Broader Vision]
\label{rem:scope}
The present paper has demonstrated \BOMA{} at the level of
foundational systems. This is the appropriate starting point, but
not the intended limit of the framework's scope.

The architectural discipline of \BOMA{} applies at every level of
mathematical development. A definition in analysis, a lemma in
algebra, a construction in topology --- each is a Block in the
sense of Definition~\ref{def:block}. A researcher who wishes to
explore two competing formulations of a definition does not face a
foundational question; they face a Decision Point at the level of
their immediate mathematical development. \BOMA{} applies there
with full generality.

This scope extension carries a methodological consequence that
the present paper has not foregrounded: the architectural
discipline makes mathematical development experimentally
tractable. Trial and error --- the exploration of alternatives
through deliberate branching, observation of consequences, and
revision of choices --- becomes a first-class mathematical
methodology rather than an informal practice. The Branch Tree is
the formal record of this experimental process; local consistency
is the criterion that makes each experimental path self-validating.

The full realization of this vision --- mathematics as a
constructible, branchable, experimentally navigable structure at
every level of development, including levels not yet reached by
existing mathematics --- is the long-term research program of
which this paper is the architectural foundation.

Within the Foundational Space $\FS{\beta_0}$, unknown mathematics
occupies a determinate structural position: the collection of
locally consistent paths that have not yet been traversed. These
paths are reachable from existing developments through finite
sequences of Decision Point reselections
(Definition~\ref{def:reachability}). The act of mathematical
exploration, within the architectural paradigm, is the act of
extending a Branch into an untraversed region of $\FS{\beta_0}$.
Mathematics is inexhaustible not because its objects are infinite
in number, but because the Foundational Space is infinite in
structure. \BOMA{} does not close this space; it makes its
infinitude navigable.
\end{remark}

\subsection{A Remark on Method}
\label{subsec:method}

This paper has been written in a style of deliberate
tentativeness. It claims what it can support and declines to
claim what it cannot. It characterizes existing foundations with
care, attributing to them the substantial achievements they have
and declining to attribute failures they do not have.

This tentativeness is not diffidence. It is an expression of the
Minimal Injection Principle applied to the domain of intellectual
claims: a paper should claim no more than is supported by its
arguments, and present no derived conclusion as an injected
assumption. The architecture is offered as the minimum structure
needed to make the architectural dimension of foundational
mathematics visible and tractable.

% ============================================================
% BIBLIOGRAPHY
% ============================================================

\newpage
\printbibliography

% ============================================================
\end{document}
% ============================================================
